\SourceLeaf{212}
\Needspace{12\baselineskip}
\begin{center}
\textit{Forms contained in others but not equivalent to them.}
\end{center}

\Needspace{8\baselineskip}
\Article{213}
Above, in Articles~157, 158, we showed that if a form $f$ with determinant $D$ contains a form $F$ with determinant $E$ and passes into it by the substitution $\alpha,\beta,\gamma,\delta$, then $E=(\alpha\delta-\beta\gamma)^2D$. If $\alpha\delta-\beta\gamma=\pm1$, the form $f$ not only contains $F$ but is equivalent to it. So if $f$ contains $F$ without being equivalent to it, the quotient $\frac{E}{D}$ is an integer greater than 1. The problem we must now solve is therefore \textit{to decide whether a given form $f$ with determinant $D$ contains a given form $F$ with determinant $Dee$}, where $e$ is assumed to be a positive number greater than 1. We will do this by showing how to give a finite collection of forms contained in $f$, such that $F$, if contained in $f$, must be equivalent to one of them.

I. Let all the positive divisors of $e$ (including 1 and $e$) be $m,m',m''$ etc., and let $e=mn=m'n'=m''n''$ etc. For brevity, denote the form into which $f$ passes by the proper substitution $m.\,0,0,n$ by $(m;0)$, the form into which $f$ passes by the proper substitution $m,1,0,n$ by $(m;1)$, etc., and generally the form into which $f$ passes by the proper substitution $m,k,0,n$ by $(m;k)$. Similarly, let $f$ pass by the proper substitution $m',0,0,n'$ into $(m';0)$, by $m',1,0,n'$ into $(m';1)$, etc., by $m'',0,0,n''$ into $(m'';0)$, and so on. All these forms are properly contained in $f$, and each has determinant $=Dee$. Denote the collection of all forms $(m;0)$, $(m;1)$, $(m;2)\ldotp\ldotp\ldotp\ldotp(m;m-1)$, $(m';0)$, $(m';1)\ldots(m';m'-1)$; $(m'';0)$ etc., whose number is $m+m'+m''+\text{etc.}$ and which are easily seen to be all distinct, by $\Omega$.

For \textit{example}, if the form $f$ is $(2,5,7)$ and $e=5$, then $\Omega$ contains the following six forms: $(1;0)$, $(5;0)$, $(5;1)$, $(5;2)$, $(5;3)$, $(5;4)$. Written out, these are $(2,25,175)$, $(50,25,7)$, $(50,35,19)$, $(50,45,35)$, $(50,55,55)$, $(50,65,79)$.

II. I now claim that if a form $F$ with determinant $Dee$ is properly contained in $f$, it must be properly equivalent to one of the forms $\Omega$. Suppose $f$ passes into $F$ by the proper substitution $\alpha,\beta,\gamma,\delta$. Then $\alpha\delta-\beta\gamma=e$. Let the greatest common divisor of $\gamma,\delta$ (which cannot both be 0), taken positively, be $=n$, and let $\frac{e}{n}=m$, which is clearly an
\SourceLeaf{213}
integer. Choose $g,h$ so that $\gamma g+\delta h=n$, and finally let $k$ be the least positive residue of $\alpha g+\beta h$ modulo $m$. Then $(m;k)$, which clearly belongs to $\Omega$, is properly equivalent to $F$, and passes into it by the proper substitution
\[
\frac{\gamma}{n}\cdot\frac{\alpha g+\beta h-k}{m}+h,\quad
\frac{\delta}{n}\cdot\frac{\alpha g+\beta h-k}{m}-g,\quad
\frac{\gamma}{n},\quad\frac{\delta}{n}
\]
For \textit{first}, these four numbers are clearly integers; \textit{second}, it is easy to check that the substitution is proper; \textit{third}, the form into which $(m;k)$ passes by this substitution is clearly the same as the one into which $f$\HistoricalNote{*)}{Since it passes by the substitution $m,k,0,n$ into $(m;k)$. See Article~159.} passes by the substitution
\[
m\left(\frac{\gamma}{n}\cdot\frac{\alpha g+\beta h-k}{m}+h\right)+\frac{k\gamma}{n},\quad
m\left(\frac{\delta}{n}\cdot\frac{\alpha g+\beta h-k}{m}-g\right)+\frac{k\delta}{n},\quad\gamma,\quad\delta
\]
or, since $mn=e=\alpha\delta-\beta\gamma$, and therefore $\beta\gamma+mn=\alpha\delta$, $\alpha\delta-mn=\beta\gamma$, by
\[
\frac{1}{n}(\alpha\gamma g+\alpha\delta h),\quad
\frac{1}{n}(\beta\gamma g+\beta\delta h),\quad\gamma,\quad\delta
\]
or finally, since $\gamma g+\delta h=n$, by $\alpha,\beta,\gamma,\delta$, \textit{i.e.}, by hypothesis, into $F$. So $(m;k)$ and $F$ are properly equivalent. \textit{This proves the claim.}

We can therefore always decide whether a given form $f$ with determinant $D$ properly contains a form $F$ with determinant $Dee$. To decide whether $f$ improperly contains $F$, we need only investigate whether the form opposite to $F$ is properly contained in $f$ (Article~159).

\Needspace{8\baselineskip}
\Article{214}
\textsc{Problem.} \textit{Given two forms, $f$ with determinant $D$ and $F$ with determinant $Dee$, the former properly containing the latter, give all proper transformations of $f$ into $F$.}

\textit{Solution.} Let $\Omega$ denote the same collection of forms as in the preceding article. Select from it all forms to which $F$ is properly equivalent; let them be $\Phi,\Phi',\Phi''$ etc. Each of these forms gives proper transformations of $f$ into $F$ in the following way, different forms giving different transformations (\textit{i.e.}, each gives distinct ones), and all together giving all of them (\textit{i.e.}, there is no proper transformation of $f$ into $F$ that is not supplied by one of the forms $\Phi,\Phi'$ etc.). Since the method is the same for all forms $\Phi,\Phi'$ etc., we will discuss just one.

\SourceLeaf{214}
Let $\Phi$ be $(M;K)$, and $e=MN$, so that $f$ passes into $\Phi$ by the proper substitution $M,K,0,N$. Further, let all proper transformations of $\Phi$ into $F$ be denoted by $\mathfrak{a},\mathfrak{b},\mathfrak{c},\mathfrak{d}$. Then clearly $f$ passes into $\Phi$ by the proper substitution $M\mathfrak{a}+K\mathfrak{c},M\mathfrak{b}+K\mathfrak{d},N\mathfrak{c},N\mathfrak{d}$. In this way every proper transformation of $\Phi$ into $F$ gives a proper transformation of $f$ into $F$. --- Treat the remaining forms $\Phi',\Phi''$ etc. in the same way: each of their proper transformations into $F$ gives a proper transformation of $f$ into $F$.

To see that this solution is complete in every respect, we must show:

I. \textit{This gives all possible proper transformations of $f$ into $F$.} Let any proper transformation of $f$ into $F$ be $\alpha,\beta,\gamma,\delta$. As in II of the preceding article, let $n$ be the greatest common divisor of $\gamma,\delta$, and determine $m,g,h,k$ as there. Then $(m;k)$ is among the forms $\Phi,\Phi'$ etc., and
\[
\frac{\gamma}{n}\cdot\frac{\alpha g+\beta h-k}{m}+h,\quad
\frac{\delta}{n}\cdot\frac{\alpha g+\beta h-k}{m}-g,\quad
\frac{\gamma}{n},\quad\frac{\delta}{n}
\]
is one of its proper transformations into $F$. From this the rule just given produces the transformation $\alpha,\beta,\gamma,\delta$. All this was proved in the preceding article.

\Needspace{26\baselineskip}
II. \textit{All transformations obtained in this way are distinct; none is obtained twice.} It is easy to see that different transformations of the same form $\Phi$ or $\Phi'$ etc. into $F$ cannot produce the same transformation of $f$ into $F$. To prove that different forms, \textit{e.g.}, $\Phi$ and $\Phi'$, also cannot give the same transformation, we argue as follows. Suppose the proper transformation $\alpha,\beta,\gamma,\delta$ of $f$ into $F$ is obtained \textit{both} from the proper transformation $\mathfrak{a},\mathfrak{b},\mathfrak{c},\mathfrak{d}$ of $\Phi$ into $F$ \textit{and} from the proper transformation $\mathfrak{a}',\mathfrak{b}',\mathfrak{c}',\mathfrak{d}'$ of $\Phi'$ into $F$. Let $\Phi=(M;K)$, $\Phi'=(M';K')$, $e=MN=M'N'$. Then we have the equations
\begin{align*}
\alpha&=M\mathfrak{a}+K\mathfrak{c}=M'\mathfrak{a}'+K'\mathfrak{c}' \tag*{. . . . . . . [1]}\\
\beta&=M\mathfrak{b}+K\mathfrak{d}=M'\mathfrak{b}'+K'\mathfrak{d}' \tag*{. . . . . . . [2]}\\
\gamma&=N\mathfrak{c}=N'\mathfrak{c}' \tag*{. . . . . . . . . . . . . [3]}
\\
\SourceLeaf{215}\delta&=N\mathfrak{d}=N'\mathfrak{d}' \tag*{. . . . . . . . . . . . . [4]}\\
\mathfrak{a}\mathfrak{d}-\mathfrak{b}\mathfrak{c}&=\mathfrak{a}'\mathfrak{d}'-\mathfrak{b}'\mathfrak{c}'=1 \tag*{. . . . . . . . . [5]}
\end{align*}
From $\mathfrak{a}[4]-\mathfrak{b}[3]$, with the help of equation [5], we get $N=N'(\mathfrak{a}\mathfrak{d}'-\mathfrak{b}\mathfrak{c}')$, so $N'$ divides $N$. Similarly, $\mathfrak{a}'[4]-\mathfrak{b}'[3]$ gives $N(\mathfrak{a}'\mathfrak{d}-\mathfrak{b}'\mathfrak{c})=N'$, so $N$ divides $N'$. Since both $N$ and $N'$ are assumed positive, we must have $N=N'$ and $M=M'$, and then [3] and [4] give $\mathfrak{c}=\mathfrak{c}'$, $\mathfrak{d}=\mathfrak{d}'$. Further, $\mathfrak{a}[2]-\mathfrak{b}[1]$ gives
\[
K=M'(\mathfrak{a}\mathfrak{b}'-\mathfrak{b}\mathfrak{a}')+K'(\mathfrak{a}\mathfrak{d}'-\mathfrak{b}\mathfrak{c}')=M(\mathfrak{a}\mathfrak{b}'-\mathfrak{b}\mathfrak{a}')+K'
\]
so $K\equiv K'\;(\mathrm{mod.}\,M)$. This can hold only if $K=K'$, since both $K$ and $K'$ lie between 0 and $M-1$. So the forms $\Phi,\Phi'$ are not distinct, contrary to the hypothesis.

Clearly, if $D$ is negative or a positive square, this method lets us actually find all proper transformations of $f$ into $F$. But if $D$ is positive and nonsquare, we can give general formulas containing all proper transformations, of which there are infinitely many.

Finally, if $F$ is improperly contained in $f$, the method given lets us easily obtain all improper transformations of the former into the latter. Namely, suppose $\alpha,\beta,\gamma,\delta$ denote all proper transformations of $f$ into the form opposite to $F$. Then all improper transformations of $f$ into $F$ are given by $\alpha,-\beta,\gamma,-\delta$.

\textit{Example.} We seek all transformations of $(2,5,7)$ into $(275,0,-1)$, which is both properly and improperly contained in the former. We gave the collection $\Omega$ for this case in the preceding article. Investigation shows that both $(5;1)$ and $(5;4)$ are properly equivalent to $(275,0,-1)$. By the theory developed above, all proper transformations of $(5;1)$, \textit{i.e.}, $(50,35,19)$, into $(275,0,-1)$ are found to be contained in the general formula
\[
16t-275u,\quad -t+16u,\quad -15t+275u,\quad t-15u
\]
where $t,u$ range over all integers satisfying $tt-275uu=1$.
\SourceLeaf{216}
So all proper transformations of $(2,5,7)$ into $(275,0,-1)$ arising from these are contained in the general formula
\[
65t-1100u,\quad -4t+65u,\quad -15t+275u,\quad t-15u
\]
Similarly, all proper transformations of $(5,4)$, \textit{i.e.}, $(50,65,79)$, into $(275,0,-1)$ are contained in the general formula
\[
14t+275u,\quad t+14u,\quad -15t-275u,\quad -t-15u
\]
and so all proper transformations of $(2,5,7)$ into $(275,0,-1)$ arising from these are contained in
\[
10t+275u,\quad t+10u,\quad -15t-275u,\quad -t-15u
\]
These two formulas therefore contain all the required proper transformations\HistoricalNote{*)}{All proper transformations are given more elegantly by the formula
\[
10t+55u,\quad t+2u,\quad -15t-55u,\quad -t-3u
\]
where $t,u$ range over all integers satisfying $tt-11uu=1$.}. --- In the same way we find that all improper transformations of $(2,5,7)$ into $(275,0,-1)$ are contained in the following two formulas:
\[
\text{(I)}\;\ldots\quad 65t-1100u,\quad 4t-65u,\quad -15t+275u,\quad -t+15u
\]
and
\[
\text{(II)}\;\ldots\quad 10t+275u,\quad -t-10u,\quad -15t-275u,\quad t+15u
\]

